\chapter{Requirement Analysis}
\label{ch:requirements}

\section{Introduction}

This chapter sets out what the platform must do. It identifies the people who use
the system, describes the main use cases with a diagram and detailed scenarios,
and then lists the functional and non-functional requirements that guided the
design. The requirements were gathered through the iterative process described in
Chapter~1, refined as each increment was built and tried on real documents.

\section{Stakeholders and Users}

The platform serves two kinds of user and depends on one external party.

\begin{itemize}[leftmargin=*,itemsep=2pt]
  \item \textbf{Member} — a signed-in user who belongs to an organisation. A
        member connects libraries, processes documents, asks questions, generates
        artefacts, explores the knowledge graph and takes part in team chat.
  \item \textbf{Team administrator} — a member with the additional ability to
        manage who belongs to the organisation and what they may access.
  \item \textbf{Model providers} — the external language and vision services
        (OpenAI and OpenRouter) that the platform calls during processing and
        answering. They are not users, but the system depends on them.
\end{itemize}

\section{Use Cases}

Figure~\ref{fig:usecase} shows the principal use cases and the actors that take
part in them. The members carry out the day-to-day work, the team administrator
additionally manages membership and access, and the external model services
participate in answering questions and generating artefacts.

\begin{figure}[H]
  \centering
  \includegraphics[width=0.96\textwidth]{fig_usecase.png}
  \caption{Use-case diagram of Synapse. A member can connect and process a
  library, ask grounded questions, generate artefacts and explore the knowledge
  graph; a team administrator additionally manages members and access. The
  external model services participate where reasoning over language or images is
  required. All use cases act within the boundary of the platform.}
  \label{fig:usecase}
\end{figure}

The most important use cases are described in detail below. Each scenario lists
the actors, the precondition, the main flow and the result.

\begin{table}[H]
  \centering
  \caption{Use case --- Connect and process a library. This is the entry point
  for all other features: until a library has been processed, there is nothing to
  question, visualise or map. The pipeline runs automatically once a source is
  connected, and the user is shown the progress of each document.}
  \label{tab:uc-ingest}
  \small
  \begin{tabularx}{\textwidth}{@{}p{0.22\textwidth}X@{}}
    \toprule
    \textbf{Use case}      & Connect and process a library \\
    \textbf{Actor}         & Member (primary); model providers (supporting) \\
    \textbf{Precondition}  & The member is signed in and belongs to an organisation. \\
    \midrule
    \textbf{Main flow}     & 1. The member connects a source of documents to a new library. \newline
                             2. The system fetches the files and detects the layout of each page. \newline
                             3. Born-digital text is read directly; scanned pages are read with a vision-language model. \newline
                             4. Figures and tables are captioned, and the text is chunked, embedded and clustered. \newline
                             5. The member sees each document move through the stages to completion. \\
    \textbf{Alternative}   & A document that fails a stage is marked with an error and can be retried without reprocessing the others. \\
    \textbf{Postcondition} & The library is processed and ready to be questioned, visualised and mapped. \\
    \bottomrule
  \end{tabularx}
\end{table}

\begin{table}[H]
  \centering
  \caption{Use case --- Ask a grounded question. This is the most frequent use of
  the platform. The reply is always grounded in the connected library and carries
  a citation for each claim, and the system abstains rather than guessing when the
  evidence is insufficient, which is what makes the answers trustworthy.}
  \label{tab:uc-chat}
  \small
  \begin{tabularx}{\textwidth}{@{}p{0.22\textwidth}X@{}}
    \toprule
    \textbf{Use case}      & Ask a grounded question \\
    \textbf{Actor}         & Member (primary); model providers (supporting) \\
    \textbf{Precondition}  & At least one processed library is selected. \\
    \midrule
    \textbf{Main flow}     & 1. The member types a question in plain language. \newline
                             2. The system analyses the question and retrieves passages by meaning and by keyword, fusing the results. \newline
                             3. It re-ranks the passages and grades whether they are sufficient. \newline
                             4. It composes an answer grounded in the passages, with a citation for each claim. \\
    \textbf{Alternative}   & If the evidence is insufficient, the system says so and does not invent an answer. \\
    \textbf{Postcondition} & The member has a cited answer, or a clear statement that the library cannot support one. \\
    \bottomrule
  \end{tabularx}
\end{table}

\begin{table}[H]
  \centering
  \caption{Use case --- Generate an artefact. The agent turns a request and a set
  of libraries into a chart, document or image. It asks a clarifying question only
  when the request is genuinely ambiguous, and the data in any chart is taken
  directly from the source files so that the figures are faithful.}
  \label{tab:uc-agent}
  \small
  \begin{tabularx}{\textwidth}{@{}p{0.22\textwidth}X@{}}
    \toprule
    \textbf{Use case}      & Generate an artefact \\
    \textbf{Actor}         & Member (primary); model providers (supporting) \\
    \textbf{Precondition}  & A library is selected, or a file is attached to the request. \\
    \midrule
    \textbf{Main flow}     & 1. The member states a goal, such as a chart or a written summary. \newline
                             2. The agent plans the work and decides what data it needs. \newline
                             3. It acquires the data by re-parsing structured files or retrieving from the library. \newline
                             4. It builds the artefact, validates it, renders it interactively and as a file, and explains the result. \\
    \textbf{Alternative}   & If the request is ambiguous, the agent asks a short clarifying question before proceeding. \\
    \textbf{Postcondition} & The member has an artefact they can view and download, with the assumptions noted. \\
    \bottomrule
  \end{tabularx}
\end{table}

\section{Functional Requirements}

Table~\ref{tab:fr} lists the functional requirements. They are grouped by the
feature they belong to, so that the requirements for ingestion, chat, the agent,
the knowledge graph and collaboration can be read together.

\begin{table}[H]
  \centering
  \caption{Functional requirements of the platform, grouped by feature. Together
  they define a system that ingests mixed documents, answers questions with cited
  and honest replies, generates faithful artefacts, maps a corpus as a graph, and
  lets teams collaborate, all within a multi-tenant boundary.}
  \label{tab:fr}
  \small
  \begin{tabularx}{\textwidth}{@{}p{0.10\textwidth}X@{}}
    \toprule
    \textbf{ID} & \textbf{Requirement} \\
    \midrule
    FR-1  & The system shall let a member connect a source and create a library of documents. \\
    FR-2  & The system shall detect page layout and read born-digital text directly. \\
    FR-3  & The system shall read scanned pages using optical character recognition and a vision-language model. \\
    FR-4  & The system shall caption figures and tables, and chunk, embed and cluster the text of each document. \\
    FR-5  & The system shall show the processing status of each document and allow a failed document to be retried. \\
    FR-6  & The system shall answer a plain-language question using only the selected libraries. \\
    FR-7  & The system shall retrieve evidence by combining vector and keyword search and shall re-rank the results. \\
    FR-8  & The system shall attach a citation to the source passage for every claim in an answer. \\
    FR-9  & The system shall abstain when the retrieved evidence is insufficient to support an answer. \\
    FR-10 & The system shall let a member request an artefact and shall ask a clarifying question only when the request is ambiguous. \\
    FR-11 & The system shall generate interactive charts and diagrams, written documents, and images, each downloadable as a file. \\
    FR-12 & The system shall take chart values directly from the source data so that artefacts are faithful. \\
    FR-13 & The system shall extract entities and relationships from a processed library and present them as an interactive knowledge graph. \\
    FR-14 & The system shall keep a link from each graph node and edge back to its source passage. \\
    FR-15 & The system shall let members of an organisation share libraries and chat together in real time. \\
    FR-16 & The system shall let a team administrator manage members and their access. \\
    FR-17 & The system shall keep each organisation's data separate from every other organisation's data. \\
  \end{tabularx}
\end{table}

\section{Non-Functional Requirements}

Table~\ref{tab:nfr} lists the non-functional requirements, which describe how
well the system must behave rather than what it must do. They cover performance,
reliability, security and privacy, scalability, usability and maintainability.

\begin{table}[H]
  \centering
  \caption{Non-functional requirements. These qualities were treated as
  first-class goals: honesty and privacy in particular shaped the design of the
  retrieval loop and the multi-tenant data model described in later chapters.}
  \label{tab:nfr}
  \small
  \begin{tabularx}{\textwidth}{@{}p{0.10\textwidth}p{0.24\textwidth}X@{}}
    \toprule
    \textbf{ID} & \textbf{Quality} & \textbf{Requirement} \\
    \midrule
    NFR-1 & Performance     & A typical question should return its first evidence within a few seconds, and a full answer within the time a user will reasonably wait. \\
    NFR-2 & Reliability     & Document processing shall be resumable, so that a failure in one stage or document does not lose the work already done. \\
    NFR-3 & Honesty         & The system shall prefer to abstain over giving an unsupported answer, and shall make its citations checkable. \\
    NFR-4 & Security \& privacy & Each organisation's documents and conversations shall be accessible only to its members, enforced at the database. \\
    NFR-5 & Scalability     & The processing pipeline shall handle large libraries by working in batches rather than all at once. \\
    NFR-6 & Usability       & A member shall be able to use every feature through the web application without technical knowledge. \\
    NFR-7 & Maintainability & The system shall be built from small, well-separated modules so that features can evolve independently. \\
  \end{tabularx}
\end{table}

\section{Assumptions and Constraints}

The platform assumes that documents are provided in common formats and that the
external model services are available, since processing and answering depend on
them. It is constrained to use established hosted and open models rather than
training new ones, and the graphics-intensive processing runs on a single
GPU-equipped server, which sets a practical bound on how much work can be in
progress at once. These assumptions match the scope set in Chapter~1.

\section{Summary}

This chapter identified the users of the platform, presented the main use cases
through a diagram and detailed scenarios, and set out seventeen functional and
seven non-functional requirements. Honesty and privacy were highlighted as
qualities that would shape the design. The next chapter turns these requirements
into a concrete design and architecture.
